% Owner: Kais
\section{Network}

\subsection{Topology}

The network has three tiers (Fig.~\ref{fig:topology}). Inside the
disconnected zone, beacons dropped by the Writers form an ESP-NOW broadcast
mesh, and every robot is a node of that mesh. At each entrance, a pair of
wired \emph{Paired Entry Beacons} (PEB) bridges the mesh to the outside.
Outside, the PEB reaches the Command Post over 4G. No robot ever talks to the
Command Post directly: every frame passes through a PEB pair.

\begin{figure}[h]
\centering
\includegraphics[width=\textwidth]{network-topology}
\caption{Network topology for one entrance.}
\label{fig:topology}
\end{figure}

\paragraph{Set-up.} Before any other task, and before taking instructions
from the Command Post, a Writer installs the infrastructure of its entrance:
PEB-Outside outside, PEB-Inside just past the entrance, joined by a
1--1.5\,m cable. Only then does it receive its GO\_TRIGGER and start exploring.

\paragraph{Data flow.} Events travel \emph{up}: Writer $\rightarrow$ beacon
mesh $\rightarrow$ PEB-Inside $\rightarrow$ cable $\rightarrow$ PEB-Outside
$\rightarrow$ 4G $\rightarrow$ Command Post. Orders travel \emph{down} the same
path. The operator chooses an Executor and sends a GO\_TRIGGER, which
PEB-Outside broadcasts over ESP-NOW to the robots waiting outside. The
briefing is split between the two halves because a briefing must happen at a
beacon: PEB-Outside only tells the Executor to \emph{enter}, and PEB-Inside
gives it the event details and its mission once it is inside.

\paragraph{Any number of robots and entrances.} Every robot has a unique
1-byte ID, so a GO\_TRIGGER can name its target Executor and every frame names
its Writer (Section~\ref{sec:message}). Each entrance has its own PEB pair. If
two meshes meet inside, a frame can leave through both pairs; the Command Post
keeps one copy per frame identifier.

\subsection{Network Communication: ESP-NOW}

All nodes inside the zone (beacons, PEBs, robots) use ESP-NOW, Espressif's
connectionless protocol on the 2.4\,GHz Wi-Fi radio:
\begin{itemize}
  \item \textbf{No infrastructure.} There is no access point, router or
    association; a node broadcasts and every node in range hears it. This
    suits a mesh whose nodes appear one by one as beacons are dropped, and
    disappear when they fail.
  \item \textbf{Fits our frames.} The 250-byte payload limit holds an event
    frame with up to 118 movement primitives ($14 + 2N$ bytes).
  \item \textbf{Low latency and low power.} A frame is sent in milliseconds,
    without a connection handshake.
  \item \textbf{One cheap chip.} The radio is built into the ESP8266/ESP32
    that each node needs anyway. Wi-Fi would need an access point, BLE has a
    shorter range, and LoRa needs an extra module and has a far lower data
    rate.
\end{itemize}
Beacons relay frames by flooding: a beacon drops a frame whose checksum is
wrong or whose \texttt{msg\_id} it has already seen, and decrements a hop
counter (TTL) before relaying. We validated relay, deduplication and TTL on a
bench of 3 ESP8266 nodes: every check passed with no send failure. Range has
not been measured yet. \todo{indoor range test}

\subsection{Paired Entry Beacons (PEB)}

The PEB is split into two halves, each built on the beacon hardware below with
its own battery, and joined by a cable carrying a UART link (TX, RX, GND,
115\,200\,baud). At 1--1.5\,m a plain UART is reliable; if the cable gets
longer or picks up motor noise, two RS-485 transceivers (about \$1 each) make
it robust over hundreds of metres. The same frames cross the cable as over the
air, each preceded by a start byte and a length byte.
\begin{itemize}
  \item \textbf{PEB-Inside} is a normal mesh beacon that also stores every
    event frame it receives, to brief Executors as they enter.
  \item \textbf{PEB-Outside} converts between the cable, ESP-NOW (to the
    robots waiting outside) and a 4G/LTE module with built-in GNSS
    \todo{module model}. The modem draws current peaks of about 2\,A, so it
    runs on a larger pack (2 $\times$ 18650).
\end{itemize}
PEB-Outside needs two serial ports, one for the cable and one for the modem.
The ESP8266 has only one full UART, so one link must use a software UART;
the ESP32 has three, which is one reason we plan to move to it.

\subsection{Beacons}

\begin{itemize}
  \item \textbf{Controller:} ESP8266 (ESP-NOW), moving to ESP32.
  \item \textbf{Power:} one 18650 Li-ion cell \todo{capacity}, its voltage
    read through a divider on the ADC.
  \item \textbf{Status LED:} off while the beacon is still carried by the
    Writer (not ready), steady on while broadcasting, blinking when the
    battery is low.
  \item \textbf{Enclosure:} 3D-printed (Section~\ref{sec:mechanics}).
\end{itemize}
A beacon keeps its frame in non-volatile memory, so it survives a reset, and
rebroadcasts it every 3\,s.

\subsection{Outside Network Area}

The Outside Network Area is made of PEB-Outside, a free cloud MQTT broker
and the Command Post. Over 4G, PEB-Outside sits behind the carrier's NAT and
cannot accept incoming connections, and neither can the Command Post laptop.
Both therefore connect \emph{out} to the broker:
\begin{itemize}
  \item PEB-Outside publishes incoming frames and its GNSS fix, and
    subscribes to GO\_TRIGGERs.
  \item The Command Post does the reverse.
\end{itemize}
Each entrance has its own topic (\texttt{livingmap/<entrance>/up},
\texttt{/down}, \texttt{/gps}), so adding an entrance needs no change at the
Command Post. The GNSS fix anchors the map of each entrance to real-world
coordinates (Section~\ref{sec:frame-translation}). The Command Post is a
laptop running our fork of the RoidOME terminal monitor, which already decodes
our frames live; reading from MQTT instead of a serial bridge is planned
for the next phase.
